\documentclass[11pt]{article}

\usepackage[margin=1in]{geometry}
\usepackage[T1]{fontenc}
\usepackage{lmodern}
\usepackage{amsmath}
\usepackage{hyperref}

\def\UrlBreaks{\do\_\do\-}

\newcommand{\Pfull}{P_{\mathrm{full}}}
\newcommand{\Psub}{P_{\mathrm{sub}}}
\newcommand{\source}[1]{\path{#1}}
\newcommand{\lean}[1]{\par\smallskip\noindent\begingroup\footnotesize\path{#1}\endgroup\par\smallskip}

\title{Theorem 2 High-Skill Individual-Fairness Direction Audit}
\author{GLM20 Dropping Standardized Testing Formalization}
\date{May 16, 2026}

\begin{document}
\maketitle

\section*{Summary}

The current source TeX supports the existing formalization caveat for
Theorem 2(ii).  The main theorem statement says that, after dropping tests,
the high-skill individual-fairness gap increases.  The appendix proof derives
the opposite high-skill comparison:
\[
  I(q;\Pfull) > I(q;\Psub)
\]
for all sufficiently high \(q\), and explicitly describes this as the gap
decreasing.  These two eventual high-skill directions cannot both be true for
the same gap definition and the same two policies.

\section*{Source Evidence}

The model defines the individual-fairness gap as
\[
  I(q;P) =
  \Pr(Y=1 \mid q,A,P) - \Pr(Y=1 \mid q,B,P),
\]
so a larger value of \(I(q;P)\) means a larger group-A-minus-group-B admission
probability gap.  This definition is in
\source{source_tex/model.tex:47--50}.

Theorem 2(ii), in the main theorem statement, says that there is a threshold
\(\hat q\) such that the individual-fairness gap increases for all
\(q>\hat q\) after dropping the test.  In context this is the comparison
\[
  I(q;\Psub) > I(q;\Pfull).
\]
This statement appears in \source{source_tex/theory.tex:184--186}.

The appendix proof of Theorem 2(ii) instead concludes the reverse comparison.
In Step 2 it sets out to discuss the high-skill direction, but the derived
sufficient condition and final conclusion are
\[
  I(q;\Pfull) > I(q;\Psub)
\]
for all sufficiently high \(q\), followed by the prose conclusion that the
gap decreases.  The relevant lines are
\source{source_tex/proofs.tex:559--596}.

\section*{Formalization Status}

Lean keeps the two directions separate.

The appendix-direction source-model endpoint is formalized by:
\lean{paper_theorem2_dropping_tests_without_barriers_standardGaussian_source_model_appendix_direction}

The main-text direction is now proved under the extra same-side precision
assumption that the remaining non-test features still favor group A:
\[
  subB < subA.
\]
The paper-facing endpoint is:
\lean{paper_interface_theorem2_high_skill_individual_fairness_gap_sub_gt_full_source_surface_of_precision_order}

This assumption matches the model/prose convention that group-B features are
less informative than group-A features, but it is not stated explicitly in the
main theorem statement.  The proof uses the reusable Gaussian left-tail
dominance lemma:
\lean{standardGaussianCDF_affine_leftTail_lt_mul_eventually_of_slope_lt}

Lean also records the direct incompatibility of the two eventual high-skill
claims:
\lean{paper_theorem2_eventual_high_skill_gap_directions_inconsistent}

\section*{Recommendation}

Keep Theorem 2 marked with a validation caveat.  If the authors intended the
main-text direction, the theorem statement should include the extra test-free
precision-order assumption \(subB < subA\), or otherwise clarify that the
``less informative for group B'' convention applies to the remaining features
after dropping the test.  Appendix Step 2 should then be corrected: its crossed
precision-order argument proves the opposite high-skill direction.  The
formalization should not mark both high-skill directions as unconditional paper
statements.

\end{document}
